\textbf{A. The analysis of the graph:}

The graph is linear:
\[
  y = ax + b = ax;
\]
From the graph it can be obtain the following data:
\begin{align*}
  \log\frac{L}{L_S} &= a\cdot\log\frac{M}{M_S}; \\
  a = \tan\alpha &= \frac{\Delta y}{\Delta x} = \frac{3{,}5}{1} = 3{,}5;
\end{align*}
% FIGURE NEEDED: the log(L/L_S) vs log(M/M_S) graph re-drawn with the slope triangle (Delta x, Delta y) and angle alpha marked (2014 TH solutions, p. 12)
\[
  L \sim M^{3,5}. \tag{1}
\]
The total energy of the star is:
\[
  E = Mc^2,
\]
So the emitted energy due to the mass variation of the star is:
\[
  \Delta E = c^2 \Delta M,
\]
According to the text
\begin{align*}
  \Delta M &= \eta M; \\
  \Delta E &= c^2 \eta M.
\end{align*}
By using the definition of the luminosity:
\begin{align*}
  \frac{\Delta E}{\Delta t} &= L; \tag{2} \\
  \Delta t &= \tau; \\
  \frac{c^2\eta M}{\tau} &= L; \\
  \tau &= \frac{c^2\eta M}{L}, \tag{2}
\end{align*}
Which represents the life-time of the star.

By using the results from the graph analysis
\[
  L = \frac{L_S}{M_S^{3,5}}\cdot M^{3,5},
\]
Thus:
\[
  \tau = \frac{c^2\eta M_S^{3,5}}{L_S}\cdot M^{-2,5}.
\]
If use the same calculations for the Sun it can be obtain
\[
  E_S = M_S c^2;
\]
\[
  \tau_S = \frac{c^2\eta_S M_S}{L_S},
\]
Which is the life-time of the Sun
\begin{align*}
  \tau &= \frac{\eta}{\eta_S}\cdot\tau_S\cdot M_S^{2,5}\cdot M^{-2,5}; \\
  \frac{\tau}{\tau_S} &= \frac{\eta}{\eta_S}\cdot\left(\frac{M}{M_S}\right)^{\!-2,5};
  \quad M = nM_S; \\
  \tau &= \frac{\eta}{\eta_S}\,(n)^{-2,5}\,\tau_S.
\end{align*}
